\section{Current Evidence Status and Results Contract}
\label{sec:results}

No quantitative research result from RQ1--RQ3 is claimed in this revision.
This is a substantive boundary rather than an unwritten section.  The archived
pilot outputs were generated under an earlier pipeline in which item pools and
the response-generating representation could vary with the evaluated KC
policy.  The current method instead fixes both before policy selection.  The
old counts therefore answer a different, confounded question and cannot be
carried forward selectively.

Table~\ref{tab:evidence-status} distinguishes what has been verified from what
still requires execution or review.  At this revision, the complete automated
suite passes 54 tests, including fixed-bank equality across ontologies,
observable/oracle separation, development-only selection, zero-KC event
retention, and non-updating probes.  These tests establish implementation
properties on controlled fixtures.  They do not estimate normalisation quality
or KT performance on the frozen EGP pilot.

\begin{table*}[t]
  \centering
  \small
  \begin{tabularx}{\textwidth}{@{}
    >{\raggedright\arraybackslash}p{.29\textwidth}
    >{\raggedright\arraybackslash}p{.27\textwidth}
    >{\raggedright\arraybackslash}X@{}}
    \toprule
    Evidence layer & Current status & Consequence for this report \\
    \midrule
    Executable specifications and boundary tests
      & Available; 54/54 tests pass
      & The declared transformations, leakage guards, frozen-probe behaviour,
        and six-policy fixed-data fixture can be described as implemented. \\
    Archived \texttt{runs/base} output
      & Predates the final pipeline; current validation fails four gates
      & Its normalisation, item, policy, and KT counts are not reported as
        results of the present method. \\
    External EGP snapshot
      & Consult-only and intentionally not redistributed
      & A fresh run requires the separately held file matching the declared
        SHA-256 and 1,222-record count. \\
    Model-assisted manual audit
      & Audit notebook checkpoints remain \textsc{not reviewed}
      & Automated validity and repeat diagnostics cannot yet be promoted to
        expert linguistic validation. \\
    Current full experimental run
      & Not yet generated and validated
      & RQ1--RQ3 quantitative estimates and ontology rankings remain pending. \\
    \bottomrule
  \end{tabularx}
  \caption{Evidence boundary at the report revision.  ``Implemented'' and
  ``empirically supported'' are deliberately treated as different statuses.}
  \label{tab:evidence-status}
\end{table*}


The archived \texttt{base} directory does not pass the current run validator.
It lacks the new \texttt{kc\_selection} stage and the required normalisation,
canonical-attrition, and repeated-item-diagnostic audit files.  The expected
external EGP file is also not distributed in the repository, and every manual
checkpoint in the research-audit notebook remains unreviewed.  Consequently:

\begin{itemize}
  \item RQ1 has an executable measurement procedure but no reportable final
        mapping yield, transition counts, or reviewed error analysis;
  \item RQ2 has a declared candidate family and selector but no validated
        selected-policy artifact from the final pilot; and
  \item RQ3 has fixture-level proof of the fixed-data and frozen-state
        boundaries but no fresh pilot event stream or ontology comparison.
\end{itemize}

Table~\ref{tab:results-contract} specifies the outputs that must replace this
status account after the first clean run.  All fields are generated from
versioned audit artifacts.  Reporting them together prevents a high predictive
score from obscuring low coverage, or a compact ontology from obscuring
unsupported distinctions.

\begin{table*}[t]
  \centering
  \small
  \begin{tabularx}{\textwidth}{@{}l X X@{}}
    \toprule
    Question & Required reported evidence & Interpretation guard \\
    \midrule
    RQ1: normalisation
      & Phase-1 and final status counts; Phase-2 transition table; admitted
        source edges and unique cells; repeat agreement by decision type;
        manually reviewed error categories and examples
      & Exact-cell yield measures conservatism and coverage, not correctness;
        automated repeats are not inter-annotator agreement. \\
    RQ2: structural support
      & Candidate support and scope diagnostics; equivalence classes; unmet
        and satisfied obligations; complete selection trace; retained KCs;
        development versus holdout activation; Q-matrix coverage, density,
        width, support, and redundancy for every condition
      & Deterministic representative choice does not distinguish equivalent
        hypotheses; a feasible inclusion-minimal cover is not a global or
        cognitive optimum. \\
    RQ3: compositional transfer
      & Identical item/event fingerprints; covered, fully supported, cold, and
        uncovered rates; all-probe log loss; secondary metrics; paired
        learner-bootstrap differences and 95\% intervals; novel-feature
        results separately
      & Comparisons are conditional on the declared synthetic oracle and
        baseline models; subset metrics cannot replace the all-probe metric. \\
    \bottomrule
  \end{tabularx}
  \caption{Result-reporting contract for the first fresh validated run.  These
  fields replace ad hoc selection of favourable counts or metrics.}
  \label{tab:results-contract}
\end{table*}


The acceptance gate for a quantitative revision is: a source digest match; a
complete run of all six ontology conditions from the same accepted bank and
event stream; \texttt{PASS} from the run validator; matching fixed-data
fingerprints in every comparison; completion of the manual normalisation and
item audits; and a result table populated directly from saved summaries.  A
failed gate remains visible and is not repaired by copying values from an
obsolete run.
